\section{Training Details}
\label{app:details}

\paragraph{SFT stage.}
We fine-tune Qwen3-32B using LoRA with rank 64 and alpha 128 on all
linear layers. The SFT dataset consists of 5{,}000 samples: 637 Chinese
mathematical critique examples and 4{,}363 English math reasoning problems
from NuminaMath-CoT and MetaMathQA. Training runs for 3 epochs (120
steps) with learning rate $5 \times 10^{-5}$, cosine decay with 3\%
warmup, batch size 2, gradient accumulation 8, and max sequence length
8{,}192. Final SFT loss: 0.311, token accuracy: 89.7\%.

\paragraph{GRPO stage.}
We train with TopoPRM rewards on 637 Chinese mathematical critique
prompts. Generation hyperparameters are configuration-dependent; in the
released runs we use temperatures in $[0.7, 0.9]$, top-$p=0.95$, and
max completion length 2{,}048, with 6--16 completions per prompt
depending on the reward-aggregation variant. We adopt the
MS-Swift \emph{server mode} architecture, where a dedicated vLLM rollout
server (4 GPUs, tensor-parallel) handles high-throughput generation while
4 separate GPUs perform LoRA gradient updates. After each optimisation
step, updated adapter weights are synchronised back to the server via
NCCL broadcast, ensuring the generation policy stays current. Training
runs for 1 epoch with learning rate $5 \times 10^{-6}$, cosine decay,
per-device batch size 2, and gradient accumulation 4.

\paragraph{Reward weights.}
For linear aggregation (\texttt{topo\_composite}), we use
$w_o = 0.4$, $w_t = 0.2$, $w_c = 0.15$, $w_f = 0.15$, $w_l = 0.1$.
For hierarchical aggregation (\texttt{topo\_hierarchical}), the base term
uses $(w_o, w_f, w_l) = (0.70, 0.15, 0.15)$ and topology/continuity
act as multiplicative gain factors (Section~\ref{sec:method}).

\paragraph{Hardware.}
All experiments run on 8$\times$ NVIDIA L20X GPUs (143\,GB VRAM each).
For GRPO, we partition GPUs into an inference group (4 GPUs running the
vLLM rollout server with tensor parallelism) and a training group
(4 GPUs running distributed LoRA optimisation), following the
MS-Swift~\citep{msswift} server-mode paradigm.

\section{SCAE Algorithm}
\label{app:scae}

Algorithm~\ref{alg:scae} presents the full procedure of Stratified
Clipping Advantage Estimation.

\begin{algorithm}[H]
\caption{Stratified Clipping Advantage Estimation (SCAE)}
\label{alg:scae}
\begin{algorithmic}[1]
\Require Group of completions $\{c^{(1)}, \ldots, c^{(G)}\}$ for prompt $q$
\Require Outcome rewards $\{R_\mathrm{out}^{(i)}\}_{i=1}^G$, auxiliary rewards $\{r_\mathrm{aux}^{(i)}\}_{i=1}^G$
\State $\bar{r}_\mathrm{acc} \leftarrow \frac{1}{G}\sum_{i=1}^G R_\mathrm{out}^{(i)}$
\State $\mathcal{C} \leftarrow \{i : R_\mathrm{out}^{(i)} = 1\}$ \Comment{correct set}
\State $\mathcal{W} \leftarrow \{j : R_\mathrm{out}^{(j)} = 0\}$ \Comment{incorrect set}
\State $\bar{r}_\mathrm{aux}^+ \leftarrow \mathrm{mean}(\{r_\mathrm{aux}^{(i)} : i \in \mathcal{C}\})$ if $\mathcal{C} \neq \emptyset$ else $0$
\State $\bar{r}_\mathrm{aux}^- \leftarrow \mathrm{mean}(\{r_\mathrm{aux}^{(j)} : j \in \mathcal{W}\})$ if $\mathcal{W} \neq \emptyset$ else $0$
\For{$i \in \mathcal{C}$}
    \State $A^{(i)} \leftarrow (1 - \bar{r}_\mathrm{acc}) + \max(0,\; r_\mathrm{aux}^{(i)} - \bar{r}_\mathrm{aux}^+)$
\EndFor
\For{$j \in \mathcal{W}$}
    \State $A^{(j)} \leftarrow -\bar{r}_\mathrm{acc} + \min(0,\; r_\mathrm{aux}^{(j)} - \bar{r}_\mathrm{aux}^-)$
\EndFor
\Return Advantages $\{A^{(i)}\}_{i=1}^G$
\end{algorithmic}
\end{algorithm}

\section{DAG Extraction Details}
\label{app:dag}

\paragraph{Extraction pipeline.}
Given a reasoning trace in \texttt{<think>}, we apply four deterministic stages:
\begin{enumerate}[leftmargin=*,itemsep=1pt]
    \item \textbf{Step segmentation}. We split the trace by numbered-step markers, mathematical discourse markers (for example ``$\because$'', ``$\therefore$''), and line-level boundaries.
    \item \textbf{Sub-question partitioning}. Markers such as ``【小题$k$】'' define sub-question identifiers. Steps from different identifiers are treated as different connected components in the final graph; cross-sub-question edges are disallowed by construction.
    \item \textbf{Information extraction}. For each step, we extract expressions, claims, and variables with rule-based patterns.
    \item \textbf{Edge construction and verification}. We add dependency edges from explicit evidence and then add weak sequential edges only when strong evidence is absent.
\end{enumerate}

\paragraph{Information-density constraints.}
To avoid low-information nodes and noisy links, we use two filters:
\begin{itemize}[leftmargin=*,itemsep=1pt]
    \item \textbf{Claim completeness}. Claims are extracted at sentence level. Heading fragments and unfinished prefixes (for example, ``all cases are:'') are removed.
    \item \textbf{Expression informativeness}. Trivial tokens (single symbols, pure numerals, and non-relational fragments) are excluded from dependency matching.
\end{itemize}
This design biases graph construction toward semantically informative units rather than formatting artifacts.

\paragraph{Dependency rules.}
For each pair of steps $(r_i, r_j)$ with $i<j$ and the same sub-question identifier, we add a \emph{virtual edge} if at least one of the following holds:
\begin{enumerate}[leftmargin=*,itemsep=1pt]
    \item \textbf{Expression reuse}: $r_j$ contains an informative expression first introduced in $r_i$.
    \item \textbf{Claim-key reuse}: canonicalized claim keys overlap between $r_i$ and $r_j$.
    \item \textbf{Variable reference}: $r_j$ reuses variables introduced in $r_i$.
\end{enumerate}
When a derivation/conclusion step has no explicit support, we insert a \emph{double-barrier edge} from the previous step within the same sub-question as a non-reward fallback connector.

\paragraph{Adaptive sequential weak edges.}
Consecutive-step weak edges (weight 0.3) are not added unconditionally.
Under the adaptive mode, a weak sequential edge is added only when the target step does not already have strong incoming support and the source step does not already emit strong outgoing dependencies.
This rule reduces chain-dominant artifacts and preserves branch structures in long traces.

\paragraph{Node typing and local verdicts.}
Each node is classified by deterministic keyword rules into
\textsc{definition}, \textsc{derivation}, \textsc{computation},
\textsc{conclusion}, \textsc{auxiliary}, \textsc{substitution}, or
\textsc{case\_analysis}.
Local verdict labels follow a hybrid strategy: reference labels are copied when available; otherwise fallback rules assign
\texttt{correct}/\texttt{incorrect}/\texttt{unverifiable} from structural support and contradiction cues.

\section{Reward Function Details}
\label{app:rewards}

\paragraph{Outcome reward.}
Parses the model's \texttt{<answer>} JSON to extract the predicted
student score and critique conclusion. Exact score match yields 1.0;
within $\pm 1$ yields 0.5. Matching the conclusion adds 0.5 bonus.
The total is capped and normalised to $[0, 1]$.

\paragraph{Format reward.}
Checks for \texttt{<think>$\cdots$</think>} and
\texttt{<answer>$\cdots$</answer>} blocks with valid JSON:
both present $\to$ 1.0; only valid \texttt{<answer>} $\to$ 0.5; only valid \texttt{<think>} $\to$ 0.1; otherwise $\to$ 0.0.

\paragraph{Length reward.}
Linear decay from 1.0 (at $\leq$ 2{,}000 characters) to 0.0 (at
$\geq$ 4{,}000 characters).

\paragraph{Topological structure reward.}
Defined in \S\ref{sec:reward_module}. The DAG extraction uses rule-based step
segmentation (numbered steps, mathematical delimiters, logical
connectives) and dependency detection (shared mathematical expressions,
variable references).

\paragraph{Continuity reward.}
Defined in \S\ref{sec:reward_module}. Uses regex-based expression matching
to verify that each step references expressions or claims from prior
steps or the problem statement.

\section{Limitations}
\label{app:limitations}

Four limitations of the current work warrant discussion. First, the DAG
extraction relies on rule-based parsing of inter-step references through
keyword matching and expression overlap. This procedure may miss implicit
logical dependencies that lack explicit surface-level cues, such as
unstated algebraic identities or geometric properties applied without
reference. Second, the continuity reward uses regex-based expression
matching rather than semantic understanding; it detects when a step
references an expression from a previous step but cannot identify subtle
logical gaps or incorrect applications of mathematical principles that
reuse the same symbols. Third, our evaluation focuses exclusively on
mathematical reasoning in a Chinese educational critique setting. Whether
TopoPRM generalises to other structured reasoning domains, such as code
generation or scientific reasoning, remains an open question.

\paragraph{Fourth limitation: TopoPRM on short adversarial benchmarks.}
On small long-CoT benchmarks (AIME~2024, AIME~2025, CNMO~2024) our
GRPO-trained TopoPRM variants (hierarchical / gated) do \emph{not}
exceed the SFT baseline on pass@1 and in some cases trail by 4--13~pp
absolute. Two concrete causes, supported by our training logs:
\begin{enumerate}[leftmargin=*,itemsep=1pt]
\item \textbf{Short GRPO budget.} The public 9B checkpoints used in the
main table are \texttt{checkpoint-79}, roughly warm-up length for GRPO.
Across the run, reward mean drifts from $0.12$ to $0.11$ without a
sustained upward trend, and \texttt{frac\_reward\_zero\_std}\,$\approx\!0.36$
throughout, meaning $\sim$40\% of each GRPO rollout group contributes
zero advantage signal. These trajectories are reported verbatim in the
supplementary file \texttt{docs/reward\_collapse\_diagnosis\_2026-04-20.md}.
\item \textbf{Tight gating threshold for the 9B gated variant.} At
\texttt{checkpoint-79} the gated reward stays at $\approx\!0.013$ with
\texttt{reward\_std}\,$\approx\!5{\times}10^{-4}$ for the entire run, indicating the
composite signal rarely crosses the gate; the gated student is thus
effectively an SFT adapter with noise, which explains why its AIME/CNMO
numbers track SFT rather than improving over it.
\end{enumerate}
The main-table narrative accommodates this: TopoPRM's claim in this
paper is structural supervision and compression, not a pass@1 ceiling
over SFT at short training budgets. Continuing GRPO to 300--500 steps
with a \texttt{reward\_std\_floor} regulariser, and annealing the gating
threshold, are left for future work
(see \S\ref{app:future} and \texttt{docs/exp\_roadmap\_2026-04-20.md}).

\section{Future Work}
\label{app:future}

Three directions follow from this work. First, replacing the rule-based
DAG extractor with a lightweight learned parser could capture implicit
dependencies while remaining fast enough for online reward computation
during GRPO training. Second, combining the verifiable structural rewards
of TopoPRM with learned PRMs~\citep{lightman2023lets} could provide
complementary supervision: structural rewards enforce global coherence
while learned rewards capture fine-grained step correctness. Third,
extending the framework to multi-modal mathematical reasoning, where
diagrams and figures introduce spatial dependency structures beyond text,
would test the generality of topology-aware rewards in richer settings.

\section{Case Study}
\label{app:case_study}

We present a representative example comparing the reasoning structure
of completions from the SFT baseline, outcome-only GRPO, and TopoPRM
on a high-school algebra problem requiring multi-step equation solving
with verification of intermediate steps.

\paragraph{SFT output.}
The SFT model produces a correct answer but with redundant intermediate
steps and one orphan conclusion that does not connect to the main
derivation chain. The reasoning DAG contains a back-edge from the
conclusion to an earlier auxiliary step, violating acyclicity.

\paragraph{Outcome-only GRPO output.}
The outcome-only model reaches the correct answer more concisely but
skips a key substitution step, producing a reasoning trace with a
continuity gap. The topological structure reward assigns a low continuity
score to this trace.

\paragraph{TopoPRM output.}
The TopoPRM model produces a concise, well-structured derivation where
each step follows logically from its predecessors. The extracted DAG is
acyclic, fully connected, and directionally consistent. The reasoning
trace is 23\% shorter than the SFT output while maintaining correctness.

\section{Distillation Details}
\label{app:distill}

For reasoning compression, we use the TopoPRM-trained Qwen3-32B as the
teacher model to generate high-quality reasoning traces. We filter
teacher outputs to retain only those with composite reward
$R_\mathrm{total} > 0.7$. The student models (Qwen3-14B, Qwen3-7B) are
then trained using reverse-KL distillation with LoRA (rank 32, alpha 64)
for 3 epochs.

\section{Additional Results}
\label{app:additional}

\paragraph{Cross-scale ablation.}
Table~\ref{tab:ablation_cross_scale} reports the effect of removing individual reward components across three model scales (Qwen3-32B, Qwen3.5-9B, Qwen2.5-7B).
The topology reward contributes the largest marginal gain at 32B scale ($\Delta$Avg = $-9.9$ when removed), while the continuity reward becomes relatively more important at smaller scales.
This suggests that larger models already produce more structurally coherent traces, making topology supervision more impactful for enforcing global DAG validity.

\paragraph{Accuracy--length trade-off.}
Table~\ref{tab:ablation_length} shows that TopoPRM achieves the best accuracy-per-token efficiency (75.5 Acc/kTok) among all GRPO variants, producing traces that are 11\% shorter than outcome-only GRPO while maintaining higher accuracy.
The length reduction is primarily driven by the topology reward eliminating redundant exploration branches and circular derivations.

\paragraph{Reward aggregation strategies.}
We compare four aggregation strategies: linear weighted sum, clipped linear, hierarchical (multiplicative topology/continuity gain), and gated (lexicographic ordering with mathematical guarantee).
The hierarchical strategy achieves the best balance between accuracy and structural quality on the light-200 protocol, while the gated strategy provides the strongest theoretical guarantee against outcome-rank reversal (Proposition~\ref{prop:gated}).
